\section{Background and Motivation}
Face recognition technology has evolved significantly in recent years, becoming an integral part of various applications ranging from security systems to user authentication in consumer electronics. The rapid advancement in deep learning techniques, particularly convolutional neural networks (CNNs), has dramatically improved the accuracy and robustness of face recognition systems.

This project is motivated by the need for a flexible, modular, and configurable face recognition system that can be adapted to various use cases while maintaining high accuracy and real-time performance. Current face recognition solutions often come as black-box implementations with limited configurability, or as research prototypes that lack robustness for real-world deployment.

\section{Problem Statement}
Implementing a complete face recognition system involves solving several interconnected challenges:
\begin{itemize}
    \item Accurate face detection in varying lighting conditions and poses
    \item Extraction of discriminative facial embeddings
    \item Efficient verification and matching against a database of known faces
    \item Real-time performance on consumer hardware
    \item Flexibility to adapt to different use cases and deployment scenarios
    \item Maintainability and configurability without code modifications
\end{itemize}

This project addresses these challenges by developing a comprehensive face recognition pipeline with modular components that can be easily configured and extended.

\section{Objectives}
The primary objectives of this project are:

\begin{enumerate}
    \item Develop a modular face recognition system with interchangeable components
    \item Implement multiple face detection methods and compare their performance
    \item Create an efficient face recognition pipeline using state-of-the-art embedding techniques
    \item Design a centralized configuration system for easy parameter tuning
    \item Support multiple operational modes (camera feed, image processing, API)
    \item Implement additional features such as face alignment and drowsiness detection
    \item Provide a web interface for easy interaction with the system
    \item Ensure the system is well-documented and maintainable
\end{enumerate}

\section{Contributions}
This project makes several contributions to the field of face recognition:

\begin{itemize}
    \item A flexible architecture that supports multiple face detection models (MediaPipe, YOLOv8, YOLOv8-ONNX)
    \item A comprehensive configuration system that centralizes all tunable parameters
    \item Integration of face alignment techniques to improve recognition accuracy
    \item Support for multiple verification metrics with configurable thresholds
    \item A modular design that allows easy extension and customization
    \item Implementation of additional features such as drowsiness detection
    \item A web-based API interface for easy integration with other systems
\end{itemize}

\section{Document Structure}
The remainder of this document is organized as follows:

Chapter 2 presents the overall system architecture, describing how the different components interact.

Chapter 3 details the configuration system, explaining how it centralizes all settings and makes the system easily adaptable.

Chapters 4-8 delve into the core components: face detection, recognition, verification, alignment, and drowsiness detection.

Chapter 9 describes the database implementation for storing and retrieving face embeddings.

Chapter 10 covers camera handling and image processing utilities.

Chapter 11 explains the API and web interface implementation.

Chapters 12-15 discuss testing, performance evaluation, deployment considerations, and potential future improvements.

The appendices provide installation guides, configuration references, API documentation, and detailed code architecture.